\cotefile{64}
\scanpage{15}
\mspage{7}

Considérons le groupe $\mathcal D$ en tant que groupe ``abstrait'', il est \unsure{canoniqt|canoniquement} extension de $\mathcal D_{\mathrm{ab}}=W$ ($\simeq\mathbb F_2\times\mathbb F_2$ non canoniqt) par $\mathbb F_2$. Dans $\mathcal D_{\mathrm{ab}}$ il y a un élément non nul distingué $\overline\omega$, \deleted{\unsure{image inverse|image réciproque}} \inserted{dont l'image} inverse est l'ens $\underline\omega_{\unsure{Q|\mathrm{le}}}$ des deux éléments d'ordre $4$ de $\mathcal D$. On voit aisément que les deux autres éléments \unsure{deux|non} nuls de $\mathcal D_{\mathrm{ab}}$ sont interchangés par le groupe des auto. de $\mathcal D$, on trouve donc
\diagram{\begin{tikzcd}[ampersand replacement=\&,column sep=small]
1\arrow[r,dash] \& \check W\arrow[r] \& \operatorname{Aut}\mathcal D\arrow[r] \& \pm1\arrow[r,dash] \& 1\\[-2pt]
\& {\substack{\text{autom de}\\\deleted{\unsure{\mathcal D}}\text{l'ext. }\mathcal D\text{ de}\\ W\text{ par }\mathbb F_2\text{ qui}\\\text{induisent l'identité}\\\text{dans }W\text{ et }\mathbb F_2}} \& \& {\underbrace{\vphantom{\pm1}\hspace{1em}}_{\substack{\text{interchange}\\\text{les les deux}\\\text{éléments de}\\\deleted{\unsure{\mathcal D_{\mathrm{ab}}}}W\setminus\{0,\overline\omega\}}}} \&
\end{tikzcd}}

% slanted
\margin{\uline{NB} on a en fait un iso. can. $\check W\simeq W$ (valable pour tt vectoriel de rg $2$ sur $\mathbb F_2$)}

En fait $\operatorname{Aut}\mathcal D$ est produit $\frac12$ direct de $\pm1$ par $\check W$, $\pm1$ opérant par échange des \inserted{$2$} éléments non \deleted{\unsure{distins|distin}} privilégiés de $\check W\simeq W$. Quand on tient compte que $\mathcal D$ provient d'un carré, donc il y a choix d'un élément $\underline c$ de $W\setminus\{0,\overline\omega\}$ (correspondant \deleted{\unsure{au}} à $C\deleted{\unsure{=|\simeq}}V_S\setminus\{1,\underline a\}$), l'autre élément étant $\underline d$ (correspondant à $D=V_A\setminus\{1,\underline a\}$). Les autom de $\mathcal D$ qui respectent cette structure \deleted{forment} le sous-groupe $\check W\simeq W$ de $\operatorname{Aut}\mathcal D$, formé des autom \inserted{d'ext.} qui induisent l'identité sur le quotient $W$. Ce sont aussi les autom intérieurs, qui forment bien le groupe $W\simeq\mathcal D/\underbrace{\{1,a\}}_{\text{centre}}$.

Comme l'hom $\mathcal D\xrightarrow{\mathrm{int}}\operatorname{Aut}\mathcal D$ n'est pas injectif, on voit que la donnée du \deleted{\unsure{carré|autre}} groupe $\mathcal D=\operatorname{Aut}Q$ (qui revient à la donnée de $C$ et de $D$, \inserted{cf plus bas}) \deleted{\unsure{comme on a} $\mathcal D=\mathbb F_2^C\mathop{\times}_{\mathbb F_2}\mathbb F_2^D$)}

% archive sheet number: 14

